\documentclass[11pt]{article}
\usepackage[margin=1in]{geometry}
\usepackage{amsmath,amssymb}
\usepackage{hyperref}
\usepackage{parskip}
\usepackage{enumitem}

\newcommand{\R}{\mathbb{R}}
\newcommand{\E}{\mathbb{E}}
\newcommand{\tauH}{\tau_{\!H}}
\newcommand{\deff}{d_{\mathrm{eff}}}
\newcommand{\Hbar}{\bar H}

\title{What We Are Doing, How, and Why \\
       \large A plain-English tour of our rate--distortion lower bound \\
       for LoRA model merging, written for a curious high-schooler}
\author{Sankalp Pathak}
\date{Companion note to \texttt{theory/theorem\_v1.tex}, written 2026-04-28}

\begin{document}
\maketitle

\begin{abstract}
This document is the friendly tour of our research project. It explains
what problem we are solving, how our main theorem works, and why the
theorem is true, using everyday language and simple pictures. We assume
only high-school algebra, the Pythagorean theorem, and a tiny bit of
probability (coin flips, averages). Every piece of math notation is
introduced in words first. If you have read the formal scaffold in
\texttt{theory/theorem\_v1.tex}, you can skim; if you have not, start
here.
\end{abstract}

\tableofcontents

%-----------------------------------------------------------------
\section{The big picture}

\subsection{What is a large language model?}

A large language model (LLM), like ChatGPT or Claude, is -- roughly --
a text autocomplete that got very good and very big. Internally, it is
a giant pile of numbers called \emph{weights}. For a modern model, the
pile has tens or hundreds of billions of numbers. When you type a
prompt, the model runs a long sequence of multiplications and additions
using those weights and spits out the next word. Train the weights well
on enough text, and the ``next-word'' machine starts to look
surprisingly smart.

For this document, we only need the following abstraction: \textbf{a
model is a very long list of numbers}. Call that list $w$. Its length
is $d$, which you should think of as ``very big.''

\subsection{What is fine-tuning?}

Suppose you have a general-purpose LLM, but you want one that is really
good at, say, answering medical questions. You take the general model,
show it a bunch of medical-question examples, and let it \emph{nudge}
its weights a little bit so it answers medical questions better. That
nudge is called \textbf{fine-tuning}.

Mathematically, the fine-tuned model is $w_0 + \tau$, where $w_0$ is
the starting (general) model and $\tau$ is the \textbf{task vector}:
the change in the weights that made it good at the new task.

\subsection{What is LoRA?}

Moving a list of $d$ numbers is expensive when $d$ has billions of
entries. \textbf{LoRA} (``Low-Rank Adaptation'') is a clever trick to
make fine-tuning cheap: instead of changing all $d$ weights, you only
change the model in a thin $r$-dimensional slice, where $r$ is much
smaller than $d$ (say $r = 8$ or $r = 64$). The slice is a
low-dimensional direction inside the high-dimensional weight space.

Analogy: if the full weight space is a huge warehouse, fine-tuning the
whole thing is like repainting every wall. LoRA is like saying ``we
only care about one corner; let's repaint just that corner.'' The
task vector $\tau$ then lives inside a tiny $r$-dimensional room
$V\subset\R^d$, not the whole warehouse.

\subsection{What is model merging, and why do we care?}

Now imagine you do LoRA fine-tuning several times, once per task:
\begin{itemize}[nosep]
  \item Task 1: medicine. Task vector $\tau_1$ in its slice $V_1$.
  \item Task 2: coding. Task vector $\tau_2$ in its slice $V_2$.
  \item Task 3: poetry. Task vector $\tau_3$ in its slice $V_3$.
  \item $\ldots$ Task $T$. Task vector $\tau_T$ in its slice $V_T$.
\end{itemize}
You now have $T$ fine-tuned models. Storing and serving $T$ of them is
expensive. So a natural question is: can you squash them into
\emph{one} merged model $w$ that is decent at all $T$ tasks at once?
This is called \textbf{model merging}.

\subsection{The rate--distortion question (our question!)}

Here is the twist. Instead of asking ``what's the best merging
algorithm'' (many people study that), we ask the opposite: \textbf{if
you only get $R$ bits to describe the merged model, what is the best
you could possibly do?} That question has a classical name:
\textbf{rate--distortion} (RD). ``Rate'' is the number of bits you
spend. ``Distortion'' is how bad your model is. An RD lower bound says
``no matter how clever you are, with $R$ bits you cannot do better than
this formula.''

Analogy. Rate--distortion is the same idea behind JPEG image
compression. A small JPEG (few bits, low rate) looks blurry (high
distortion); a big JPEG (many bits, high rate) looks sharp (low
distortion). There is a fundamental trade-off between ``how few bits''
and ``how good it looks,'' and the RD lower bound tells you the
best-case curve.

We prove such a lower bound, but for \emph{merged LLM task vectors}
instead of images. \textbf{This is the central contribution of our
paper.}

%-----------------------------------------------------------------
\section{The exact question we are answering}

\subsection{Setup in one picture}

You have $T$ arrows $\tau_1,\dots,\tau_T$, all starting at the origin
of a $d$-dimensional space. Each arrow has length $B$ (we fix a common
length for normalization; $B$ is just a constant, think $B=1$). Each
arrow lives inside its own $r$-dimensional slice $V_t$.

You must pick a single arrow $w$ (the merged model) that does the job
of all of them. But you are only allowed to describe $w$ using $R$
bits.

\subsection{How do we measure ``did a good job''?}

Two natural options.
\begin{itemize}[nosep]
  \item \textbf{Average error:} $\frac{1}{T}\sum_{t=1}^T \|w - \tau_t\|^2$.
        (The average squared distance from $w$ to each arrow.)
  \item \textbf{Worst-case error:} $\max_{t=1,\dots,T}\|w - \tau_t\|^2$.
        (The squared distance to the arrow $w$ is farthest from.)
\end{itemize}

We use \textbf{worst-case}, because that is the honest measure of a
merged model. If $w$ is great at 9 tasks out of 10 but terrible at the
10th, you probably do not want to ship it: the whole point of merging
is that it should be decent at \emph{all} tasks. Also, average error
collapses to a boring problem (just quantize the mean arrow $\bar\tau
= T^{-1}\sum_t \tau_t$), while worst-case error is genuinely new.

\subsection{A little refinement: $H_t$-weighted error}

There is one more detail. Not every direction in $\R^d$ actually
matters for task $t$. Only directions \emph{inside the slice $V_t$}
do -- moving $w$ perpendicular to $V_t$ cannot change task $t$'s loss,
because LoRA only moved things in $V_t$ in the first place. Formally,
we use the $H_t$-weighted error
\[
  (w - \tau_t)^\top H_t (w - \tau_t),
\]
where $H_t$ is an operator that ``keeps only the $V_t$ part'' of an
arrow (the orthogonal projection onto $V_t$). Think of $H_t$ as a
flashlight that only lights up directions task $t$ cares about. The
error only counts what the flashlight sees.

\subsection{The RD function}

Putting it together, define
\[
  \mathcal{D}^\star(T,d,r,B,R)
  \;=\; \text{smallest possible worst-task error},
\]
where the minimum is over all possible encoders and decoders (all ways
of compressing to $R$ bits and decompressing back to a merged $w$), and
where the task vectors $(\tau_t, V_t)$ are drawn from a specific random
family we will describe in \S\ref{sec:setup}. We prove a
\textbf{lower bound} on $\mathcal{D}^\star$: a formula that no merging
scheme can beat.

%-----------------------------------------------------------------
\section{Math setup, gently}\label{sec:setup}

Before stating our main theorem, we need five concepts.

\subsection{$d$-dimensional space}

An arrow in 2D is a pair of numbers $(x_1, x_2)$. In 3D, it is a
triple $(x_1, x_2, x_3)$. In $d$ dimensions, it is a list of $d$
numbers. The length (norm) of the arrow $\tau$ is
$\|\tau\| = \sqrt{x_1^2 + x_2^2 + \cdots + x_d^2}$
(Pythagoras, generalized). A \textbf{subspace} of dimension $r$ is a
flat slice through the origin that looks like $\R^r$ sitting inside
$\R^d$. In 3D, 2-dim subspaces are planes through the origin, and
1-dim subspaces are lines through the origin.

\subsection{The subspace $V_t$ of task $t$}

Task $t$'s LoRA fine-tuning moves weights in an $r$-dim slice $V_t$.
Its task vector $\tau_t$ therefore sits inside $V_t$, with length $B$.
That is our only structural assumption. We will average over
\emph{random} choices of $V_t$ (``Stiefel-random,'' which is math-speak
for ``pick an $r$-dim slice uniformly at random''), independently for
each task.

\subsection{The projection $P_{V_t}$ (and why $H_t = P_{V_t}$)}

$P_{V_t}$ is the flashlight. Given any arrow $x\in\R^d$, $P_{V_t}x$ is
the part of $x$ that lies inside $V_t$. The part outside $V_t$ is
thrown away. By construction, $P_{V_t}\tau_t = \tau_t$ (since $\tau_t$
is already in $V_t$).

We set $H_t = P_{V_t}$. Why? Because the actual task-$t$ loss, near
$\tau_t$, looks locally like $(w - \tau_t)^\top H_t (w - \tau_t)$ where
$H_t$ is the ``Hessian'' (curvature matrix) of the loss. For LoRA,
this $H_t$ is exactly supported on $V_t$ (directions outside $V_t$
don't affect the loss), and in the cleanest approximation,
$H_t = P_{V_t}$. It is both mathematically convenient and physically
meaningful.

\subsection{The centroid $\tauH$}

We need one ``summary arrow'' of the $T$ task vectors, chosen to
minimize the $H_t$-weighted errors. Define the average of the
flashlights:
\[
  \Hbar \;:=\; \frac{1}{T}\sum_{t=1}^T H_t
  \;=\; \frac{1}{T}\sum_{t=1}^T P_{V_t}.
\]
Then set
\[
  \tauH \;:=\; \Hbar^+\cdot\Bigl(\frac{1}{T}\sum_t H_t\tau_t\Bigr),
\]
where $\Hbar^+$ is the Moore--Penrose pseudoinverse (the closest
thing to $1/\Hbar$ that makes sense; you can ignore the ``pseudo'' part
for the high-level picture). The point of $\tauH$: it is the unique
arrow such that, when you sit at $\tauH$ and look around through each
task's flashlight, you see a nice balance.

A very useful identity: because $P_{V_t}\tau_t = \tau_t$, a short
calculation gives
\[
  \Hbar\,\tauH \;=\; \bar\tau \;:=\; \tfrac{1}{T}\sum_t \tau_t.
\]
So $\tauH$ is, in a precise sense, the ``inverse-warped mean.''

\subsection{The effective dimension $\deff$}

Here is the single most important new quantity in our paper:
\[
  \deff \;:=\; \text{rank}\Bigl(\sum_{t=1}^T P_{V_t}\Bigr).
\]
In words: the total number of independent directions that \emph{at
least one} task cares about. Two cartoons.
\begin{itemize}[nosep]
  \item \textbf{All tasks share the same slice} ($V_1 = V_2 = \cdots =
        V_T$). Then the union is just that one slice, of dimension
        $r$. So $\deff = r$.
  \item \textbf{All tasks have fully orthogonal slices.} Then the
        union has dimension $Tr$ (each slice contributes its own $r$
        fresh dimensions). So $\deff = Tr$.
\end{itemize}
The general case is somewhere between $r$ and $Tr$ (and always at most
$d$). Classroom analogy: if every student writes on the same wall,
$\deff$ is the size of that wall. If every student has their own
private wall, $\deff$ is the sum of all walls. The $\deff$ parameter
is exactly how the theorem measures ``how much the task slices
overlap.''

%-----------------------------------------------------------------
\section{The main result (Theorem 7), in plain words}\label{sec:thm}

Here is our main theorem (Theorem~7 in \texttt{theorem\_v1.tex}):

\begin{quote}
\textbf{Main Theorem.} For merged LoRA models with $T$ tasks, each
living in a random $r$-dim slice of $\R^d$, at rate $R$ bits, the
worst-case $H_t$-weighted error satisfies
\begin{equation}\label{eq:main}
  \mathcal{D}^\star
  \;\geq\;
  \underbrace{B^2\!\left(1 - \frac{\deff}{Tr}\right)}_{\text{(A) the floor}}
  \;+\;
  \underbrace{c_{\mathrm{TQ}}\cdot\frac{B^2\deff}{Tr}\cdot 2^{-2R/\deff}}_{\text{(B) the compression term}}.
\end{equation}
\end{quote}

There are two pieces, each with its own story.

\subsection{(A) The floor: what you can't fix with more bits}

The first term, $B^2(1 - \deff/(Tr))$, does not depend on $R$ at all.
It is the error you are stuck with \emph{even if $R = \infty$}, i.e.\
even if you are allowed infinite bits. It is a geometric fact about
how the task vectors point: because they sit in different slices, no
single arrow $w$ can perfectly match all of them at once. The floor
captures that unavoidable geometric gap.

Two sanity checks:
\begin{itemize}[nosep]
  \item \textbf{Shared subspace:} $V_t$ all the same, $\deff = r$.
        Floor $= B^2(1 - r/(Tr)) = B^2(1 - 1/T)$. This matches the
        classical centroid-error formula: even in the shared-slice
        case, you cannot perfectly represent $T$ different unit arrows
        by one arrow; the unavoidable loss is $B^2(1 - 1/T)$.
  \item \textbf{Orthogonal subspaces:} $\deff = Tr$. Floor
        $= B^2(1 - 1) = 0$. \textbf{Zero!} Does that make sense? Yes:
        if the tasks live in completely non-overlapping slices, then
        inside task $t$'s flashlight, no other task is visible, and
        you can stitch a single $w$ that matches each $\tau_t$ exactly
        on its own flashlight. There is no unavoidable floor.
\end{itemize}
(This second bullet surprised us -- we had initially expected the
floor in the orthogonal case to be \emph{large}, not zero. See
\S\ref{sec:numerics} for the story.)

\subsection{(B) The compression term: how bits help}

The second term, $c_{\mathrm{TQ}}\cdot(B^2\deff/(Tr))\cdot 2^{-2R/\deff}$,
is the part that shrinks as you spend more bits. The factor $2^{-2R/\deff}$
says: each extra bit per effective dimension cuts the squared error by
a factor of 4. This is classical rate--distortion scaling. Intuition:
quantizing an arrow in a $\deff$-dimensional space takes about $\deff$
bits per unit of precision; at $R$ bits total you get $R/\deff$ bits
per dimension; squared-error shrinks like $4^{-R/\deff}$. The
$c_{\mathrm{TQ}}$ is a universal constant (the constant of the
TurboQuant lower bound of Zandieh et al., arXiv 2504.19874), which we
import from existing work.

The multiplier in front, $B^2\deff/(Tr)$, is the ``compressible
budget'': the amount of the total energy that can still be improved by
more bits.

\subsection{The conservation law: a pretty observation}

Here is something beautiful that falls out. The floor and the
compression-term prefactor add up:
\[
  \underbrace{B^2\!\left(1 - \frac{\deff}{Tr}\right)}_{\text{floor}}
  \;+\;
  \underbrace{B^2\frac{\deff}{Tr}}_{\text{compression prefactor}}
  \;=\; B^2.
\]
So the total squared-length budget ($B^2$ per task) splits into
``floor (cannot fix with bits)'' plus ``compression prefactor (can fix
with enough bits),'' in a ratio controlled exactly by
$\deff / (Tr)$. When $\deff = Tr$ (orthogonal case), the entire budget
is compressible -- no floor. When $\deff = r$ (shared case), only a
$1/T$ fraction is compressible -- the rest is floor. The single
parameter $\deff / (Tr) \in [1/T, 1]$ controls the whole trade-off.

\subsection{Three endpoint checks}

The theorem recovers known limits at three boundaries (each is a
one-line substitution; see \texttt{theorem\_v1.tex} \S6).
\begin{enumerate}[nosep]
  \item \textbf{Full rank} ($r=d$, i.e.\ LoRA's slice is the full
        space): the bound collapses to our Phase 0 Theorem 1, which is
        the classical multi-task result.
  \item \textbf{Single task} ($T=1$): the bound collapses to
        TurboQuant Theorem 3, the state-of-the-art single-task RD lower
        bound.
  \item \textbf{All tasks share one slice} ($V_t = V$): the bound
        matches Phase 0 with $d \mapsto r$, exactly as you would hope.
\end{enumerate}
So our theorem properly generalizes three prior results, unified into
one formula parameterized by $\deff$.

%-----------------------------------------------------------------
\section{Why the theorem is true (proof sketch in plain English)}

The actual proof in \texttt{theorem\_v1.tex} \S5 is six steps. Here is
the same skeleton, one paragraph per step, with no hairy algebra.

\paragraph{Step 1 (Yao's minimax).} To prove a lower bound on
$\mathcal{D}^\star$, we specify a \emph{hard} distribution $P$ on
task-vector tuples and argue: no merging scheme, however cleverly
designed, can beat a certain error against this $P$. This is Yao's
minimax principle. Analogy: to show ``no student can ace this class,''
we pick one extra-tough exam and show that every study strategy fails
on it. Our exam is: each $V_t$ is a uniformly random $r$-dim slice,
and $\tau_t$ is a uniformly random length-$B$ arrow inside it.

\paragraph{Step 2 (max $\geq$ average).} The worst-case error is at
least the $H_t$-weighted average error (obviously -- the biggest of
several numbers is at least their average). The weighted average has
a nice form: a floor term (\S\ref{sec:thm}-A) plus $(w - \tauH)^\top
\Hbar (w - \tauH)$. So we have replaced the hard-to-reason-about
$\max$ with a single quadratic error against $\tauH$, plus an additive
floor we already computed in Lemma 6.

\paragraph{Step 3 (restrict to relevant directions).} Only the part of
$w$ inside $\range(\Hbar)$ -- the union of all task slices -- actually
affects the error. Components of $w$ orthogonal to every $V_t$ are
invisible to every flashlight, so WLOG we assume $w$ lives entirely in
this $\deff$-dim subspace. Analogy: if flashlights only light up a
wall, what's behind you is irrelevant.

\paragraph{Step 4 (Euclidean change of variable).} The quadratic form
$(w-\tauH)^\top\Hbar(w-\tauH)$ is not quite the familiar squared
distance, because of the $\Hbar$ weighting. But we can stretch
coordinates by $\Hbar^{1/2}$ (the square root of $\Hbar$), and in the
stretched coordinates, the quadratic form \emph{becomes} ordinary
squared distance. Specifically: let $\eta := \Hbar^{1/2}\tauH$ and
$\xi := \Hbar^{1/2}(w-\tauH)$. Then our quadratic becomes $\|\xi\|^2$
and the problem reduces to: \emph{compress $\eta$ using $R$ bits, in a
$\deff$-dim Euclidean space.} A perfectly standard
rate--distortion problem, just in a smaller-than-$d$ dimension.

\paragraph{Step 5 (symmetry $\to$ known bound).} Now the magic.
Because the slices $V_t$ were drawn uniformly at random, the resulting
$\eta$ has a highly symmetric distribution: it is rotationally
invariant inside its subspace. (Rotate the subspace any way you like,
$\eta$ looks statistically the same.) Plugging this into a known
rate-distortion lemma for rotationally-invariant sources (our Phase 0
Lemma 2, which in turn builds on TurboQuant), we get
\[
  \text{expected squared error}
  \;\geq\;
  c_{\mathrm{TQ}}\cdot\E\|\eta\|^2\cdot 2^{-2R/\deff}.
\]
This is the $2^{-2R/\deff}$ decay in the main theorem, dropping out
cleanly.

\paragraph{Step 6 (combine).} We need $\E\|\eta\|^2$. But
$\|\eta\|^2 = \tauH^\top\Hbar\tauH = B^2 - (\text{floor})$, by a short
algebra step. Plugging $\E\|\eta\|^2 = B^2\deff/(Tr)$ (from Lemma 6)
into Step 5 and combining with Step 2's floor gives exactly
\eqref{eq:main}. Done.

%-----------------------------------------------------------------
\section{Numerical experiments: did it actually work?}\label{sec:numerics}

Proofs alone are not enough -- we also write code to \emph{simulate}
the setup and check that the empirical numbers match the formula. Two
Python files do this:
\texttt{code/synthetic/day8\_rank\_r\_sanity.py} and
\texttt{code/synthetic/day9\_verify\_lemma6.py}.

\subsection{What we ran}

For each $(T, d, r)$ combination:
\begin{enumerate}[nosep]
  \item Sample random task slices $V_t$ and random task vectors
        $\tau_t \in V_t$ of length $B$ (thousands of trials).
  \item Compute the centroid $\tauH$, the effective dimension $\deff$,
        and the floor.
  \item Quantize $\tauH$ using various bit budgets $R$; measure the
        squared error as a function of $R$.
  \item Compare the empirical error-vs-bits curve to
        \eqref{eq:main}.
\end{enumerate}

\subsection{Three checks we wanted to pass}

\paragraph{(a) Regression at $r=d$.} At the full-rank boundary we
should recover Phase 0. We did: empirical floor / Phase-0 floor is in
$[0.998, 1.017]$ across all cells. Check.

\paragraph{(b) Slope of log error vs bits should be $-2$.} The
$2^{-2R/\deff}$ decay predicts that a log-log plot of excess error
against bits should have slope $-2$ per extra bit per dim. Originally
we measured this using the wrong metric (Euclidean squared error of
$w - \tauH$ in the raw basis) and saw slope $-1.2$. This was initially
alarming! After some head-scratching, we realized: the theorem
predicts the slope on the \textbf{$\Hbar$-weighted} error, i.e.\ on
$\|\eta\|^2$, not on the raw Euclidean error. When we switched to the
right metric (the one Step 4 of the proof uses), the slope became
$-2.00\pm 0.04$ across every cell we tested. Check.

Lesson learned: bugs can live in measurement, not in math. The proof
was correct; our first measurement was wrong.

\paragraph{(c) $\deff$ matches the rank formula.} We can fit the
measured $\deff$ from the error-vs-bits curve and compare to
$\rank(\sum_t P_{V_t})$. They match within 5\% across $r \in
\{4, 8, 16, 32\}$ at $T = 2$ with $Tr < d$. Check.

\subsection{A special sanity test for Lemma 6}

We also wrote \texttt{day9\_verify\_lemma6.py} that directly checks the
closed-form floor formula $\E[\text{floor}] = B^2(1 - \E[\deff]/(Tr))$.
Across 6 cases (shared $V$, orthogonal $V$, mixed; different $T, r$),
the empirical floor matched the theoretical floor to relative error
$\leq 4\times 10^{-3}$. In the orthogonal case it matched to
$10^{-24}$, i.e.\ numerical precision.

%-----------------------------------------------------------------
\section{The bug story (honest debrief)}\label{sec:bugs}

Science is not a straight line. Two real bugs that we caught during
Day 9:

\subsection{The Step-5 conditioning bug}

In our first draft of the proof, Step 5 started with ``condition on
$V_1, \dots, V_T$.'' This is wrong! If you condition on where each
individual slice sits, you have fixed the problem and there is no
longer any randomness in the $V_t$'s to average over -- the symmetry
argument collapses. The correct move is to condition \emph{only} on
the union $W = V_1 + \cdots + V_T$ (the $\deff$-dim subspace), and
keep the $V_t$'s inside $W$ random. Under that weaker conditioning,
rotational symmetry inside $W$ is preserved, and the argument goes
through. We caught this by re-reading the proof carefully \emph{after}
our numerics confirmed rotational invariance (otherwise we might have
missed it and shipped a subtle gap).

\subsection{The Marchenko--Pastur confound}

When we first ran the rotational-invariance sanity test, we saw a
spread of $1.27$ in the eigenvalues of the sample covariance of
$\eta$, which naively suggests ``the distribution is not rotationally
symmetric -- there's a preferred direction.'' But this was finite-
sample noise! Random matrix theory (the Marchenko--Pastur
distribution) predicts a spread of exactly this size even for a truly
isotropic source at this sample size. We added an MP-based upper bound
on the expected spread to the sanity test; the observed spread is
below this upper bound; symmetry confirmed. The takeaway:
\emph{statistical tests need their own theory} -- you cannot just
``eyeball the number.''

Both of these are in the daily log and discussed more in
\texttt{notes/daily\_log.md} (Day 9 entry).

%-----------------------------------------------------------------
\section{What comes next}

Our work so far covers:
\begin{itemize}[nosep]
  \item The closed-form floor (Lemma 6).
  \item The full lower bound Theorem 7 for $H_t = P_{V_t}$ (the
        LoRA-projector case).
  \item The extension to general $H_t \succeq 0$ with arbitrary
        per-task curvature spectra $(D_t)_{t=1}^T$
        (Theorem 8, added Day 10--11). This is the
        \textbf{MSE $\to$ CE bridge} -- real LLM cross-entropy
        loss has a Fisher-information Hessian that is rank-$r$ PSD
        but not a projector, and Theorem 8 applies to it directly.
        The key trick was finding the ``right'' hard distribution:
        sample $\tau_t$ uniformly on the $H_t$-ellipsoid in $V_t$
        (not the Euclidean sphere), which makes the $H_t$-weighted
        energy isotropic per-task regardless of the eigenvalue
        spectrum.
  \item Matching numerical verification for all of the above
        (Day 8, 9, 10, 11 Python scripts).
\end{itemize}

The remaining steps on the way to ICLR 2027 (deadline 2026-09-24):
\begin{itemize}[nosep]
  \item \textbf{Achievability (Phase 2, starting now).} A lower
        bound says ``no one can do better than X.'' We also want a
        matching \emph{upper bound}: a concrete merging algorithm
        (based on Hadamard incoherence + TurboQuant-style scalar
        quantization of $\eta = \bar H^{1/2}\tauH$) that achieves
        the lower bound up to constants. Then the RD function is
        pinned down: $\mathcal{D}^\star \asymp \text{floor} +
        B^2(\deff/Tr)\cdot 2^{-2R/\deff}$.
  \item \textbf{Empirical validation.} On real LLM-with-LoRA data
        (not just synthetic Stiefel-random slices), do the slopes
        still match? This is the experimental part of the paper.
\end{itemize}

If both land by summer 2026, we have a clean story for ICLR 2027:
a matching lower and upper bound on the fundamental compressibility
limit of merged LoRA models, verified both on synthetic and real data.

%-----------------------------------------------------------------
\section*{A note on the audience}

This document is written to be readable at about a 9th-grade level,
but the underlying math -- despite the analogies -- is very real.
Everything here has a rigorous version in \texttt{theorem\_v1.tex}, and
every claim was numerically verified before being written down. We
hope this document is useful as an on-ramp for readers who want to
understand what the research is about before diving into the formal
scaffold.

\end{document}
